% =====================================================================
%  SAMPLE LAYOUT of the final report. It is NOT the report.
%
%  Real content (keep):   title, section headings, equations, figures,
%                         tables, reference list.
%  Yellow boxes (note):   a paragraph YOU write goes here. The box says
%                         what the paragraph must do. Delete the box
%                         when the paragraph is written.
%  Blue boxes (model):    example sentences about a DIFFERENT topic with
%                         invented numbers. Delete all of them.
%
%  To see the bare layout without any box: change \notestrue to
%  \notesfalse on the line marked  <<< SWITCH  and compile again.
%
%  Compile:  tectonic report_sample.tex
%  Fonts: see the three cases in the comment of the font block below.
% =====================================================================
\documentclass{article}

\PassOptionsToPackage{numbers,sort&compress}{natbib}
\usepackage[preprint]{neurips_2024}
\makeatletter
\renewcommand{\@noticestring}{AA228/CS238 final project, Stanford University, Autumn 2026.}
\makeatother

\usepackage{amsmath,amssymb}
\usepackage{booktabs}
\usepackage{graphicx}
\usepackage{xcolor}
\usepackage{enumitem}
\usepackage{textcomp}
\usepackage{url}

% --- font block ---------------------------------------------------------
% (a) While the Chinese notes are in the file: keep the three commands as
%     they are and compile with tectonic or XeLaTeX on a Mac (PingFang SC
%     is a macOS font).
% (b) Notes gone, still tectonic or XeLaTeX (also on Overleaf): replace
%     the two xeCJK lines by \usepackage{fontspec} and KEEP \setmainfont.
%     Without \setmainfont these engines silently set the text in Latin
%     Modern instead of Times, and the tables get wider.
% (c) Notes gone, pdfLaTeX: delete all three commands; the style file
%     then supplies Times by itself.
\usepackage{xeCJK}
\setCJKmainfont{PingFang SC}
\setmainfont{texgyretermes}[Extension=.otf,UprightFont=*-regular,BoldFont=*-bold,
  ItalicFont=*-italic,BoldItalicFont=*-bolditalic,Ligatures=TeX]
% --- end of the font block ------------------------------------------------

\usepackage{environ}
\usepackage[most]{tcolorbox}

\newif\ifnotes
\notestrue          % <<< SWITCH: \notestrue shows the boxes, \notesfalse hides them

\definecolor{noteback}{RGB}{255,249,219}
\definecolor{noteframe}{RGB}{150,105,0}
\definecolor{modelback}{RGB}{232,241,252}
\definecolor{modelframe}{RGB}{38,84,150}
\tcbset{sample/.style={breakable,enhanced,boxrule=0.4pt,arc=1.5pt,
  left=5pt,right=5pt,top=3pt,bottom=3pt,before skip=5pt,after skip=6pt,
  fontupper=\small}}

% \begin{note}[which paragraph, how long] ... \end{note}
\NewEnviron{note}[1][]{\ifnotes
  \begin{tcolorbox}[sample,colback=noteback,colframe=noteframe]
  {\bfseries\color{noteframe}你来写\ifx\relax#1\relax\else:#1\fi}\par\smallskip
  \BODY
  \end{tcolorbox}\fi}

% \begin{model} ... \end{model}
\NewEnviron{model}{\ifnotes
  \begin{tcolorbox}[sample,colback=modelback,colframe=modelframe]
  {\bfseries\color{modelframe}示范句(别的题目:路口红绿灯;数字是编的;只学写法)}\par\smallskip
  \BODY
  \end{tcolorbox}\fi}

\newenvironment{pts}{\begin{itemize}[nosep,leftmargin=1.15em,topsep=1pt]}{\end{itemize}}

% caption placeholder: the caption is yours to write
\newcommand{\capnote}[1]{\ifnotes{\color{noteframe}【标题你来写】#1}\else Caption to be written.\fi}

% figures and tables are generated by code; show a frame if a file is missing
\newcommand{\fig}[2]{\IfFileExists{#1}{\includegraphics[width=#2]{#1}}%
  {\fbox{\parbox[c][1.5in][c]{0.9\linewidth}{\centering figure file not built yet:\\ \texttt{\detokenize{#1}}}}}}
\newcommand{\tab}[1]{\IfFileExists{#1}{\input{#1}}%
  {\fbox{\parbox[c][0.8in][c]{0.9\linewidth}{\centering table file not built yet:\\ \texttt{\detokenize{#1}}}}}}

\newcommand{\qprev}{q_{\mathrm{prev}}}
\newcommand{\rest}{\mathcal{J}}          % the set of forecasts not yet bought
\DeclareMathOperator{\E}{\mathbb{E}}

\title{2\textcent{} for Your Two Cents:\\ When Is an Unseen Forecast Worth Buying?}

% Use the names on the course roster and fill in the e-mail addresses.
\author{%
  Leonard Collomb\\ Stanford University\\ \texttt{(email)}
  \And
  Henry Chen\\ Stanford University\\ \texttt{(email)}}

\begin{document}

% =====================================================================
%  PAGE 0: how to read this sample (disappears with \notesfalse)
% =====================================================================
\ifnotes
\pagestyle{empty}
\begingroup\small
\begin{center}
{\Large\bfseries 样稿说明}\\[3pt]
{\small 这两页说明不是论文的一部分。论文从它们后面那一页开始。\\
英文版是 \texttt{report\_sample\_en.pdf},内容相同。}
\end{center}

\subsection*{这是什么}
\begin{pts}
\item 一份排好版的样稿。版式、章节、公式、图、表和参考文献都是真的;正文段落是空的,留给你们写。课程规定报告正文必须是你们自己的话,结论必须是你们自己的判断,所以这里没有替你们写任何一段英文正文。
\item 英文要你们自己直接写。不要把提示或中文草稿交给翻译软件或 AI 去生成英文,那样就不是自己的话。
\item \textbf{提示是给你们的指令,不是报告的文字。}每一句都要自己写。不要把提示翻译一下、或者改几个词就放进报告。
\item 提示里有时写了一种读法(例如``损失不是贪心带来的'')。那只是一种可能的读法:先对照图表自己判断,同意才用。
\end{pts}

\subsection*{三种内容}
\begin{note}[某一段,约多少词]
黄色框:这里要有一段你们自己写的英文。框里写的是这段话的作用、要用的数字、要引的课本位置。写好一段,就把这个框整个删掉。
\end{note}
\begin{model}
Blue box: example sentences in English. They are about a different topic, and the numbers are invented. Look at how the sentences are built. Do not copy them, and do not reuse them with the nouns changed. Delete every blue box before you submit.
\end{model}
没有底色的内容是真实的,可以留下:标题、章节名、公式、图、表、参考文献。
\begin{pts}
\item 公式是从代码里照实写出来的。你们要自己手推一遍,能讲清每一项。记号可以改。
\item 图和表由代码生成,数字不要手改。图表的标题(caption)要自己写。
\end{pts}

\subsection*{怎么用}
\begin{enumerate}[nosep,leftmargin=1.5em]
\item 先通读一遍,看清每节多长、图表放在哪。
\item 写作顺序:第 2、3 节和 4.1 节 $\rightarrow$ 4.2、4.3 节 $\rightarrow$ 第 5、6 节 $\rightarrow$ 第 1 节 $\rightarrow$ 摘要。
\item 想看不带提示的纯版面:把源文件里的 \texttt{\textbackslash notestrue} 改成 \texttt{\textbackslash notesfalse},重新编译。这时参考文献列表是空的,BibTeX 会给一个警告,因为引用都写在提示框里;这是正常的。
\item 写完发给我(Claude Code,做这份样稿的 AI 工具)。我核对数字、逻辑和课本引用,只指出问题,不替你们改写。
\item 提交:PDF;一个组员在 Gradescope 上交;截止时间是 12 月 4 日(周五)下午 5 点(太平洋时间)。作者名用选课名单上的名字,邮箱自己填。
\item 交之前检查:把 \texttt{\textbackslash notestrue} 改成 \texttt{\textbackslash notesfalse};按源文件里字体那一段注释的 (b) 或 (c) 改字体设置;在 PDF 里搜索中文字符和 ``Caption to be written'',都不应该出现;邮箱填好;数一下页数。
\end{enumerate}

\subsection*{提示里用到的词}
这几行是给你们解释词的,不是报告的文字,不要抄进报告。
\begin{pts}
\item \textbf{agent}:每天做决定的那个程序。它不训练预测模型,只看它们的输出。
\item \textbf{free forecast}($y_0$)、\textbf{forecast A}($y_A$)、\textbf{forecast B}($y_B$):对 $r$(SPY 下一个交易日的价格变化)的三个预测。免费预测每天都能看到,不花钱。A 和 B 要买了才能看到(每个 0.1 bp,每个只能买一次)。
\item \textbf{position}($q$):agent 持有到下一个交易日的仓位。\textbf{long}($q=+1$):价格涨就赚。\textbf{flat}($q=0$):什么都不持有。\textbf{short}($q=-1$):价格跌就赚。
\item \textbf{bp}(basis point):1 bp $= 0.01\% = 10^{-4}$。所有的量都是本金的比例,用 bp 报告。
\item \textbf{net utility}:一天的得分。式 (1) 减去当天买预测付的价格。式 (1) 有三部分:赚或亏,减去风险惩罚,减去换仓成本。
\item \textbf{belief}:在已经看到的预测之下,agent 对 $r$(以及还没看到的预测)的概率分布。\textbf{belief mean}($m$):这个分布里 $r$ 的均值。
\item \textbf{policy}:决定每天买哪些预测的规则。五个策略:
  \begin{itemize}[nosep,leftmargin=1.15em,topsep=0pt]
  \item \textbf{never buy}:从不花钱买 A 或 B。它每天仍然选仓位,用的规则和其他策略相同,只是只看免费预测。
  \item \textbf{always buy}:每天先买 A,再买 B。
  \item \textbf{random}:随机地买(规则见 4.1 节的第 3 个提示)。
  \item \textbf{greedy}:只往前看一次购买。净值(式 (4))最大的那个预测,净值为正就买;否则停止。
  \item \textbf{exact}:把当天的两次购买决定放在一起规划(式 (5))。
  \end{itemize}
\item \textbf{always long}、\textbf{always flat}:表 1 里的两个参考行。每天都做多(或每天都空仓),什么都不买。
\item $\boldsymbol{\lambda}$:式 (1) 里风险惩罚的权重。$\lambda$ 越大越怕风险。实验用 $\lambda=1$ 和 $\lambda=5$。
\item \textbf{interval}:95\% 的 block bootstrap 区间。它是围绕一个平均值的范围,表示这个平均值有多不确定;做法是成块地重抽样日子。区间包含 0 时,提示里说``没有结论''。
\item \textbf{fit period}:1999-01-26 到 2014-12-31(4{,}010 天),belief 在这些天上拟合。\textbf{test period}:2015-01-02 到 2026-11-12(2{,}984 天),没有用来拟合 belief;策略在这些天上比较。
\item \textbf{predicted gain}:agent 的模型认为``照 exact 做''比``立刻停止(不买)''多得多少。\textbf{realized gain}:同一个差,实际的结果。\textbf{same starting position}:每一天,两边都从 never buy 前一天的仓位出发。完整的定义见 4.1 节的第 4 个提示。
\end{pts}

\subsection*{篇幅(课程要求 4--6 页;哪些部分计入,要问助教)}
\begin{center}\small\renewcommand{\arraystretch}{0.88}
\begin{tabular}{@{}lll@{}}
\toprule
节 & 词数 & 里面有什么\\
\midrule
摘要 & 150 & 问题、设定、方法、结果、结论\\
1 Introduction & 420 & 课本的两句话、两个问题、已有工作\\
2 Problem Formulation & 220 & 图 1、式 (1)\\
3 Approach & 420 & 式 (2)--(5)\\
4.1 Setup and metrics & 360 & 数据、模型、策略、两个度量\\
4.2 Greedy and exact & 230 & 表 1、图 2\\
4.3 Predicted and realized & 230 & 表 2、表 3\\
5 Discussion and Limitations & 350 & 局限、未来工作\\
6 Conclusion & 70 & 不出现新内容\\
五个图表的标题 & 约 200 & 每个 30--50 词\\
\bottomrule
\end{tabular}
\end{center}
\begin{pts}
\item 合计约 2{,}450 词,另加标题约 200 词。用占位文字按这些词数填满后实测:正文(到 Conclusion 结束)约 5.8 页;加上 Contributions 和 AI 声明约 6.2 页;参考文献另占约 1 页。这次测试里 AI 声明约 130 词,Contributions 几乎没有字。这些页数有大约 $\pm 0.3$ 页的出入;真实文字里公式多的话会更长。
\item 所以没有余量。这周就问助教:参考文献、Contributions 和 AI 声明算不算在 6 页里。
\item 如果 Contributions 和 AI 声明也算:把表 3 换成两句话(省约 0.3 页),实测约 5.9 页。如果连参考文献也算:要压到约 1{,}700 词,不要表 3,标题每个约 20 词。
\end{pts}

\subsection*{图表怎么重新生成}
在 \texttt{cs238/} 目录里运行下面两条命令:\\
{\footnotesize\texttt{python3 -m experiments.report\_figures --no-notes --out ../final\_report/figures}\\
\texttt{python3 -m experiments.report\_tables\ \ --no-notes --out ../final\_report/tables}}\\[2pt]
{\raggedright \texttt{--no-notes}:表格下面不带英文注释行,说明由你们写进 caption。结果数字都在 \texttt{tables/numbers\_sheet.csv} 里(图里的数字在 \texttt{figures/figure\_numbers.csv}),带来源,正文里的数字从那里抄。表里 \texttt{meaning} 那一列是代码的说明文字:只抄数字,不要抄那一列的英文。配置参数(模型设置、日期)见 \texttt{cs238/configs/default.toml} 和 \texttt{cs238/results/tables/exp1\_lam1\_manifest.json}。有三类数字不在这些文件里,用到它们的提示里写了出处:数据下载日期、输入特征的个数、10 月模拟器的数字。\par}
\endgroup
\clearpage
\pagestyle{plain}
\setcounter{page}{1}
\fi

% =====================================================================
%  THE PAPER
% =====================================================================
\maketitle
\suppressfloats[t]   % no figure above the title on the first page

\begin{abstract}
\begin{note}[摘要,约 150 词,全文写完以后再写]
五句话,一句只做一件事:
\begin{pts}
\item 问题:课本 6.6 节有两句话,你们对每一句各问一个问题。第一个问题课本自己提了;第二个是你们自己的问题,课本那句话没有覆盖这种情况。
\item 设定:每天交易前,agent 可以花钱买两个还没看到的预测。
\item 方法:贪心做法和精确解都算出来,在没有用来拟合 belief 的 2{,}984 天上比较。
\item 结果,数字从这里选:两种做法在 2{,}984 天里每天选的仓位都相同;$\lambda=1$ 时模型预测每天多得 $+0.84$ bp,实际是 $-2.70$ bp(95\% 区间 $[-4.6,\,-1.0]$);$\lambda=5$ 时区间包含 0,没有结论。
\item 结论:这两个结果合起来说明什么,你们自己判断。亏损的原因没有检验过,不写。
\end{pts}
\end{note}
\begin{model}
We ask whether a traffic signal that reacts to queue lengths leaves fewer cars waiting through a full cycle than a fixed timer does. In 30 simulated evenings per controller, 11.5\% of cars waited a full cycle under the reactive signal and 18.2\% under the timer; the spread across runs was about 2 points for each. We simulated a single intersection, so the result says nothing yet about a network.
\end{model}
\end{abstract}

% ---------------------------------------------------------------------
\section{Introduction}
% ---------------------------------------------------------------------

\begin{note}[第 1 段,约 70 词:从课本出发]
\begin{pts}
\item 课本 6.6 节给了一个做法:每次挑信息价值最大的变量去观察,直到不值得再观察为止 \citep{kochenderfer2022algorithms}。
\item 这一节里有两句话,可以加引号和页码直接引用:\\
第 121 页:``The value of information is never negative.''\\
第 122 页:``This greedy selection of observations is only a heuristic; it may not represent the truly optimal sequence of observations.''
\end{pts}
\end{note}

\begin{note}[第 2 段,约 80 词:我们的问题]
\begin{pts}
\item 一天里发生什么,一两句话,指向图 \ref{fig:day}(a)。
\item 用同学熟悉的例子类比:课本 Exercise 20.5(第 424 页:依次做诊断检查,再决定治疗还是回家)。检查对应预测,治疗与否对应仓位。
\item 和那个例子的不同:那里检查可以重复做,这里每个预测只能买一次;那里各次检查在已知病情时相互独立,这里即使已知 $r$,两个预测仍然相关。还可以提:这里的量是连续的,并且有一个预测是免费的。
\item 第一次出现时用一句话定义:SPY、position(long 做多 / flat 空仓,即 $q=0$ / short 做空)、bp。
\end{pts}
\end{note}

\begin{note}[第 3 段,约 70 词:两个问题]
\begin{pts}
\item 问题一:贪心做法和精确解选的一样吗?
\item 问题二:agent 的模型预测``买了会更好'',这个好处在后来的数据上出现了吗?
\item 用一句话把标题和这两个问题连起来。标题是双关(花钱买别人的意见);实验里的价格是每个预测 0.1 bp。
\item 如实说明:问题是看到结果以后定下来的,比 proposal 窄。proposal 计划把损失拆成 planning、estimation、misspecification 三部分;这里只检验了第一部分,后两部分没有分开。
\end{pts}
\end{note}

\begin{note}[第 4 段,约 50 词:答案预告,以及不声称什么]
\begin{pts}
\item 问题一:两种做法每天选的仓位相同。问题二:$\lambda=1$ 时好处没有出现,实际反而更差;$\lambda=5$ 时没有结论。
\item ``never negative'' 说的是 agent 自己的模型下、扣掉价格之前的期望效用。你们检验的是这个保证范围之外的情况,不是反驳课本。这句话必须有。
\item 不解释亏损的原因。
\item 这一节里 $r$、$q$ 和 $\lambda$ 都还没有定义(第 2 节才定义)。第 2 段和这一段要么用文字说(next-day price change;风险惩罚的权重较小、较大的两种设定),要么加半句指向第 2 节。
\end{pts}
\end{note}

\paragraph{Related work.}\ifnotes\mbox{}\par\vspace{-3pt}\fi
\begin{note}[约 150 词:已有工作。评分表有 10 分看这一段和参考文献]
四组文献,每组一两句话说它做了什么;最后一句说你们哪里不同。
\begin{pts}
\item 信息价值和``只看一步''的规则:\citet{howard1966information}(概念的来源);\citet{russell1991principles};\citet{frazier2008kg}(belief 相互独立时的一步规则,叫 knowledge gradient);\citet{frazier2010paradoxes}(信息的边际价值不是凹的,所以一步规则可能停得太早);\citet{hay2012selecting}(定义了 myopic policy,指出它容易停得太早;它的 Theorem 7:myopic 决定继续收集时,最优策略也会继续)。
\item 贪心和最优的比较:\citet{krause2009optimal} 在图模型里直接比较了两者,并证明信息价值不是 submodular 的;\citet{golovin2011adaptive} 的贪心保证需要 adaptive submodularity,所以不覆盖你们的目标函数。
\item 来源或 belief 互相相关时:\citet{frazier2009correlated};\citet{clemen1985limits}(来源互相依赖时,信息的价值下降得很快)。
\item 预测组合、估计误差、事后失望:\citet{bates1969combination};\citet{smith2006optimizer};\citet{welch2008comprehensive}(注意:它研究的是月度到五年的预测,不是日度)。
\item 你们的不同:不要说``比较贪心和精确解''是新的,\citet{krause2009optimal} 和 \citet{hay2012selecting} 都比较过。不同在设定(为一个交易仓位购买预测、高斯 belief、有换仓成本、每个预测只能买一次)和问题二(模型预测的好处 对 后来数据上的实际好处)。
\item 只引你们读过的。这一段的论文里,项目文件夹里没有全文的有五篇:Howard 1966、Frazier and Powell 2010、Bates and Granger 1969、Clemen and Winkler 1985、Smith and Winkler 2006。先通过学校图书馆拿到再引。另外有七条引用也没有本地全文(Hoerl and Kennard 1970、Breiman 2001、Friedman 2001、Pedregosa et al. 2011、Harris et al. 2020、Virtanen et al. 2020、Politis and Romano 1992),它们是作为方法和软件的出处来引的。Politis and Romano 1992 的页码范围来自二手来源,交之前核对一次。
\end{pts}
\end{note}

% ---------------------------------------------------------------------
\section{Problem Formulation}
% ---------------------------------------------------------------------

\begin{figure}[t]
\centering
\fig{figures/fig1_one_day.pdf}{\linewidth}
\caption{\capnote{一句说图里是什么:(a) 一天里的决策,用 2015-09-10 这个真实的一天,$\lambda=1$;(b) 昨天做多时,三种仓位的期望效用随 belief 均值的变化。一句说读者该看到什么:买了两个预测以后,均值从 $+2.12$ 降到 $-2.49$ bp,跌破了``继续做多''的阈值 $+0.64$ bp,所以转为空仓。必须写明这一天是你们挑的例子:它是买了两个预测的 53 天之一;典型的一天只买 A(2{,}816 天)。还要说明:空心圆是买之前的 belief,实心圆是买之后的,两个面板里是同一对点;(b) 的曲线已经扣了换仓成本。}}
\label{fig:day}
\end{figure}

\begin{note}[第 1 段,约 40 词:补上定义。时间线在 Introduction 讲过,这里不重复]
\begin{pts}
\item $r$ 的定义,只定义一次:从今天收盘到下一个交易日收盘,经分红调整的收盘价(adjusted close)的相对变化。
\item 给 $r$ 起名叫 next-day price change。课本定义的 return(第 134 页)是奖励之和,换一个词可以避免混淆。
\item 说清 agent 是一个程序,它不训练预测模型,只看它们的输出。
\end{pts}
\end{note}

\begin{note}[第 2 段,约 100 词:问题的形式化。下面两种写法选一种,不要混着写]
\begin{pts}
\item 写法一,POMDP(像 Exercise 20.5 的解答那样列):隐藏的状态是 $r$ 和还没看到的预测;状态里看得到的部分是昨天的仓位、已经买了哪些预测、免费预测 $y_0$;动作是 buy A、buy B、stop and choose a position;观察是买到的那个预测的值;奖励是买的时候付出价格、停止时得到式 \eqref{eq:utility}。最后加一句:按课本 20.1 节把它当作 belief-state MDP 来解。
\item 写法二,belief-state MDP(课本 20.1 节,第 407 页):状态是 belief 的均值和方差、已经买了哪些预测、昨天的仓位;动作同上;转移是买到预测以后按式 \eqref{eq:cond} 更新 belief;奖励是买的时候 $-c_j$、停止时是式 \eqref{eq:stop} 的期望效用(课本式 (20.1):belief 下奖励的期望)。
\item 两种写法都要写:每天最多两次购买决策,加一次选仓位;天与天之间只通过昨天的仓位相连;规划器不跨天往前看。最后这一条是局限,第 5 节还要再提。
\end{pts}
\end{note}

\begin{equation}
u(q, r) \;=\; q\,r \;-\; \lambda\,q^{2} r^{2} \;-\; \tau\,\lvert q - \qprev \rvert ,
\qquad q \in \{-1, 0, +1\}.
\label{eq:utility}
\end{equation}

\begin{note}[第 3 段,约 80 词:逐项解释式 (1)]
\begin{pts}
\item 第一项是赚或亏。第二项是风险惩罚:$\lambda$ 越大越怕风险。第三项是换仓成本,$\tau = 1$ bp。
\item 和课本的联系:前两项就是课本式 (6.3) 的二次效用 $U(x)=\lambda x - x^{2}$,自变量是当天的盈亏 $x = q\,r$。关系要写准确:$x - \lambda x^{2} = \lambda\,[(1/\lambda)\,x - x^{2}]$,也就是课本的形式取参数 $1/\lambda$,再乘以 $\lambda$。这个系数在这里不能省,因为 $\tau$ 和价格用的是同一个单位。这一步是自己推出来的,课本没有这样写,所以用 ``our $\lambda$'' 这类说法。
\item 定义 net utility:式 \eqref{eq:utility} 减去当天买预测付的价格。第 4 节的表和图用的都是它。
\item 单位:所有量都是本金的比例,用 bp 报告(1 bp $= 10^{-4}$)。
\end{pts}
\end{note}

\begin{model}
The state is the pair of queue lengths, each between 0 and 20 cars. Every 5 s the controller takes one of two actions: keep the current green light or switch it. The reward is minus the number of waiting cars, so a larger reward means less waiting.
\end{model}

% ---------------------------------------------------------------------
\section{Approach}
% ---------------------------------------------------------------------

\subsection{Belief and its update}

\begin{note}[约 80 词]
\begin{pts}
\item $(r, y_0, y_A, y_B)$ 的联合高斯,在前 4{,}010 天上估计一次,之后不再更新。
\item 怎么估计的:均值和方差是样本值;相关系数用 Ledoit--Wolf 的收缩强度(0.023)向 0 收缩 \citep{ledoit2004wellconditioned}。所以写 ``estimated'',不要写 maximum likelihood。第一次用 shrinkage 这个词时解释它。
\item 每看到一个预测,用式 \eqref{eq:cond} 更新。这是课本 3.9 节的式 (3.25)--(3.26)(第 63 页),要引。
\item 一个值得写的性质:更新后的方差不依赖看到的数值。
\item 式 \eqref{eq:cond} 前面先用一句话定义记号:$o$ 是已经看到的预测;$\mu_r$、$\mu_o$ 是拟合出的均值;$\Sigma_{rr}$ 是 $r$ 的方差,$\Sigma_{oo}$ 是 $o$ 的协方差,$\Sigma_{ro}$ 是两者之间的协方差(课本记作 $A$、$B$、$C$;这里换了字母,因为 A、B 已经用来叫两个预测)。
\item belief $b$ 是给定 $o$ 以后,$r$ 和还没看到的预测的联合分布;其中 $r$ 的部分是 $\mathcal{N}(m, v)$。
\end{pts}
\end{note}

\begin{equation}
m \;=\; \mu_r + \Sigma_{ro} \Sigma_{oo}^{-1} (o - \mu_o), \qquad
v \;=\; \Sigma_{rr} - \Sigma_{ro} \Sigma_{oo}^{-1} \Sigma_{ro}^{\top}.
\label{eq:cond}
\end{equation}

\subsection{Choosing a position}

\begin{equation}
\E_b\!\left[u(q, r)\right] \;=\; q\,m - \lambda\,q^{2}\,(m^{2} + v) - \tau\,\lvert q - \qprev\rvert ,
\qquad
G(b) \;=\; \max_{q}\; \E_b\!\left[u(q, r)\right].
\label{eq:stop}
\end{equation}

\begin{note}[约 80 词]
\begin{pts}
\item 这是课本 6.4 节的最大期望效用原则。$G(b)$ 是``现在停下来''的价值。
\item 阈值取决于昨天的仓位,两个都要给。昨天空仓:$m - \lambda(m^{2}+v) > \tau$ 才做多,大约是 $m > \lambda v + \tau$。昨天做多:只要 $m$ 大约高于 $\lambda v - \tau$ 就继续做多。
\item 算术写出来:$r$ 的标准差约 128 bp,所以 $v \approx 1.64\times10^{-4}$,$\lambda v \approx 1.64$ bp($\lambda=1$)。两个阈值是 2.64 bp 和 0.64 bp;$\lambda=5$ 时约 9.2 bp 和 7.2 bp。
\item 昨天做多或空仓时,均值落在两个阈值之间(宽度 $2\tau$),agent 就保持昨天的仓位。图 \ref{fig:day}(b) 画的是 $\lambda=1$、昨天做多的情况。
\item 这两个阈值解释了 $\lambda$ 的作用,第 5 节还要用到。
\end{pts}
\end{note}

\subsection{Value of one more forecast}

\begin{equation}
\mathrm{VOI}_j(b) \;=\; \E_{y \sim b}\!\left[\,G\!\left(b^{j,y}\right)\right] - G(b),
\qquad
\mathrm{NV}_j(b) \;=\; \mathrm{VOI}_j(b) - c_j .
\label{eq:voi}
\end{equation}

\begin{note}[约 60 词]
\begin{pts}
\item 先用一句话定义记号:$j$ 是一个还没买的预测,$c_j$ 是它的价格,$b^{j,y}$ 是看到它的数值 $y$ 之后的 belief,$\mathrm{NV}_j$ 是买它的净值。
\item 这是课本式 (6.9)(第 121 页),把 $EU^{*}$ 换成 $G$,求和换成期望。
\item $\mathrm{VOI}_j(b) \ge 0$ 就是课本那句 ``never negative'':在 agent 自己的模型下,扣价格之前。
\item 高斯模型下式 \eqref{eq:voi} 的期望有闭式($G$ 是 $m$ 的分段二次函数),不需要采样。一句话即可。
\end{pts}
\end{note}

\subsection{Greedy rule and exact solution}

\begin{note}[约 50 词:贪心做法]
\begin{pts}
\item 对每个还没买的预测算式 \eqref{eq:voi} 的净值 $\mathrm{NV}_j$;最大的净值为正就买它,否则停止并选仓位。
\item 这就是课本 6.6 节最后一段的做法。用课本的名字称呼它(greedy),并指明是 6.6 节。文献里叫 myopic,提一次就够。
\end{pts}
\end{note}

\begin{equation}
V(b, \rest) \;=\; \max\Bigl\{\, G(b),\;\; \max_{j \in \rest}\,\bigl[\, -c_j + \E_{y \sim b}\, V\bigl(b^{j,y},\, \rest \setminus \{j\}\bigr) \bigr] \Bigr\},
\qquad V(b, \varnothing) = G(b).
\label{eq:bellman}
\end{equation}

\begin{note}[约 80 词:精确解]
\begin{pts}
\item 先定义记号:$\rest$ 是还没买的预测的集合。
\item 式 \eqref{eq:bellman} 是对 belief 的 forward search:课本 22.2 节式 (22.1)(第 454 页),折扣为 1,求和换成积分,``停止''是终止动作。
\item 深度 2 就覆盖了一整天:买完两个预测以后只剩``停止''这一个决策,所以叶子上的 $G$ 是精确的,搜索没有被截断。
\item 怎么算的:最后一步的期望是式 \eqref{eq:voi} 的闭式;式 \eqref{eq:bellman} 里第一步的期望用确定性的数值积分,没有随机采样。积分用的是复合 Gauss--Legendre 公式:把标准化以后的积分范围 $[-12, 12]$ 在被积函数的折点处切开,再切成宽度不超过 1 个标准差的小段,每段 201 个节点,总共至少 4{,}824 个节点。配置文件里的 \texttt{n\_quad = 201} 是每一小段的节点数。所以 ``exact'' 的意思是:除了这个数值积分以外是精确的。
\item 两种做法只可能在第一个决策上不同:买了一个预测以后,对剩下那个预测,贪心的判断和精确解的判断是同一个判断。这一点要写,第 5 节要用。
\item 写明它是 agent 自己模型下的精确解。不要只写 ``optimal'':它只在 agent 的模型里是最优的。
\end{pts}
\end{note}

\begin{note}[约 70 词:为什么用这个方法。评分表 Appropriateness of Approach 的 10 分在这一段]
\begin{pts}
\item 预测在买之前看不到,所以这是带信息收集动作的序贯问题。课本第 122 页说,最优的观察顺序可以用后面章节的序贯决策方法求出来;你们用的正是它。
\item 只有两个预测,精确解算得出来,才能和贪心直接比。
\item 高斯模型让更新和一步的信息价值都有闭式,所以两种做法的差别里没有采样误差。
\item 评估方式也各用半句话说明理由:为什么基准是 never buy(同一个模型、同一套选仓位的规则,只是不买,所以差别只来自购买);为什么在没有用来拟合的日子上比(问题二问的就是后来的数据)。
\end{pts}
\end{note}

% ---------------------------------------------------------------------
\section{Experiments}
% ---------------------------------------------------------------------

\subsection{Setup and metrics}

\begin{note}[第 1 段,约 60 词:数据]
\begin{pts}
\item SPY 的每日数据,用经分红调整的收盘价。来源:Yahoo Finance,用 yfinance 包下载,下载日期 2026-11-16(UTC;这个日期的出处是 \texttt{cs238/data/raw/SPY\_daily\_manifest\_ci.json})。数据范围 1993-01-29 到 2026-11-13。
\item 用来拟合 belief 的时期:1999-01-26 到 2014-12-31,4{,}010 天。测试期:2015-01-02 到 2026-11-12,2{,}984 天。测试期的数据没有用来拟合 belief。
\item 第一次用 held-out 这个词时解释它。
\end{pts}
\end{note}

\begin{note}[第 2 段,约 80 词:三个预测模型和软件。评分表要求引用不是你们自己开发的算法和软件]
\begin{pts}
\item 免费的 $y_0$:ridge regression \citep{hoerl1970ridge}。A:random forest \citep{breiman2001random}。B:gradient boosting \citep{friedman2001greedy}。都用 scikit-learn \citep{pedregosa2011scikit}。
\item 每 21 个交易日重新训练一次,只用当时已有的数据;至少 1{,}260 天训练数据。
\item 主要超参数各半句(ridge 的 $\alpha=10$;随机森林 300 棵树、深度 6;梯度提升 200 棵树、深度 2、学习率 0.05)。输入特征写一句:16 个(7 个不同窗口的对数收益、3 个已实现波动率、另外 6 个;出处是 \texttt{cs238/src/data\_pipeline.py} 里的 \texttt{FEATURE\_COLUMNS})。
\item 其余软件:数值计算用 NumPy \citep{harris2020numpy} 和 SciPy \citep{virtanen2020scipy}。
\item 用一句话给出代码位置:仓库地址、commit、配置文件。其余设置读者可以在那里查到。引用之前,先把还没有提交的东西提交:3 个脚本、4 个结果表(CSV)和 3 个测试文件。另外,回放日志和预测存档不在仓库里(被 .gitignore 排除),所以别人克隆以后不能直接运行生成图表的两个脚本。用一句话说明这些文件怎么得到(例如:由仓库里的 GitHub Actions 工作流重新生成;写之前自己确认一次)。
\end{pts}
\end{note}

\begin{note}[第 3 段,约 80 词:参数和比较的策略]
\begin{pts}
\item 参数:每个预测 0.1 bp;$\tau = 1$ bp;$\lambda \in \{1, 5\}$。
\item 价格扫描:另外用 8 个价格各跑了一遍(0、0.01、0.03、0.1、0.3、1、3、10 bp)。4.2 节要用,所以这里要先交代。
\item 策略,每个半句:never buy(基准)、always buy(每天先 A 后 B)、random(每一步以 0.5 的概率再买一个,从还没买的里面随机挑;只模拟了一次)、greedy、exact。
\item fixed order(先 A 后 B,只决定何时停)两种处理选一种:不提;或者提一句``它每天的购买和仓位都和 exact 完全相同,所以表 \ref{tab:policies} 里没有它''。
\item 两个参考行:always long、always flat。
\item 写明:每个策略带着自己昨天的仓位进入下一天。
\end{pts}
\end{note}

\begin{note}[第 4 段,约 140 词:两个度量。每个只定义一次。问题二全靠度量二,要写足]
\begin{pts}
\item 度量一:某个策略和 never buy 在同一天的 net utility 之差,取平均,单位 bp/天,带 95\% 区间。
\item 区间用 circular block bootstrap \citep{politis1992circular}:每块 21 天,2{,}000 次重抽样。课本没讲,用一句话解释为什么按块抽(相邻的日子互相有关)。
\item 稳健性检查:5 种块长(5、10、21、42、63 天)$\times$ 20 个 bootstrap 种子,共 100 个区间。
\item 区间端点只报一位小数,和表里一样。原因:端点会随 bootstrap 的随机种子变,连第一位小数也会动零点几 bp(例如 $\lambda=1$ 的 exact,20 个种子:下端从 $-6.2$ 到 $-5.9$,上端从 $-1.1$ 到 $-0.9$)。所以``区间不含 0''这句话靠的是 100 个区间的检查,不是表里印的那一位小数。
\item 度量二,用论文里的记号定义:某一天的 predicted gain 是 $V(b,\{A,B\}) - G(b)$,也就是 agent 的模型认为``照精确解做''比``立刻停止''多得多少;realized gain 是这一天照精确解做实际得到的 net utility,减去立刻停止实际得到的。两边每天都从 never buy 的昨日仓位出发。
\item 一句话说明它为什么是一个检验:如果拟合的模型是对的,realized gain 的平均应该等于 predicted gain 的平均(只差抽样误差)。
\item 一句话说明表 \ref{tab:policies} 和表 \ref{tab:gain} 的数字为什么不同:表 \ref{tab:policies} 里每个策略沿着自己的仓位路径走;表 \ref{tab:gain} 里两边每天都从同一个仓位重新出发。
\item 读数规则写在这里:区间包含 0 就说``没有结论'',不说``没有效果''。
\end{pts}
\end{note}

\begin{model}
We count a car as delayed when it waits longer than one full signal cycle. For each controller we run 30 simulations and report the share of delayed cars, averaged over the runs, together with its standard deviation across runs.
\end{model}

\subsection{Greedy and exact purchasing}

\begin{table}[t]
\caption{\capnote{一句说表里是什么(每个策略的平均 net utility,以及和 never buy 的差);一句说怎么读区间;一句说明 greedy 和 exact:每天的仓位都相同,两行只差在 1 天和 8 天多付的 0.1 bp。表格本身不带注释,这些说明都要写在这里。}}
\label{tab:policies}
\centering\small
\tab{tables/tab1_policies.tex}
\end{table}

\begin{note}[第 1 段,约 70 词:直接回答问题一]
\begin{pts}
\item 两个 $\lambda$ 下,greedy 和 exact 在 2{,}984 天里每一天选的仓位都相同。
\item 购买不同的日子:$\lambda=1$ 时 1 天,$\lambda=5$ 时 8 天。这 9 天全是 greedy 停止、exact 买了 A,没有反过来的;exact 在这些天预测的优势都不到 0.001 bp。
\item 这个方向和 \citet{hay2012selecting} 的 Theorem 7 一致(他们的设定不同,写 ``in a related setting'')。
\item 所以表 \ref{tab:policies} 里这两行不是两个独立的证据,是同一个比较。
\end{pts}
\end{note}

\begin{note}[第 2 段,约 80 词:读表 1]
\begin{pts}
\item $\lambda=1$:exact 比 never buy 每天低 3.36 bp,区间 $[-6.0,\,-1.0]$;always buy 低 3.80 bp,区间 $[-6.6,\,-1.2]$。
\item random 低 1.35 bp,区间 $[-2.7,\,0.0]$。这是边缘结果:区间上端其实是 $-0.005$,表里印成 0.0;100 个稳健性区间里 74 个不含 0;random 本身只模拟了一次。
\item $\lambda=5$:所有差的区间都包含 0,没有结论。
\item 两件必须点明的事:never buy 和 always long 分不开(4.42 对 4.40,差 $-0.02$,区间 $[-2.0,\,+1.9]$);$\lambda=5$ 时 always flat 的点估计最高。
\item 表里的差是用没有舍入的数算的,所以 $3.06-4.42$ 和表里的 $-1.35$ 差 0.01 是舍入,不是错误。
\end{pts}
\end{note}

\begin{figure}[t]
\centering
\fig{figures/fig2_cumulative.pdf}{\linewidth}
\caption{\capnote{一句说纵轴是什么(相对 never buy 的累计 net utility 之差,bp);一句说读者该看到什么(差距是在哪些年份累积的;两个面板用同一个纵轴)。说明累计值是每天差值的简单相加,不是复利。线末端的陡降是真实数据:最后 7 个交易日里 exact 相对 never buy 少了约 754 bp。}}
\label{fig:cumulative}
\end{figure}

\begin{note}[第 3 段,约 80 词:两种做法在别的价格下也一致;差距何时累积]
\begin{pts}
\item 价格扫描:8 个价格、两个 $\lambda$,共 16 次运行,greedy 和 exact 每天的仓位都相同。
\item 这不是因为第二个预测很少被买:$\lambda=1$、价格为 0 时,89.7\% 的日子两个预测都买了,仓位仍然相同。(价格 0.1 bp 时是 1.8\%;0.3 bp 及以上是 0。)
\item 图 \ref{fig:cumulative}:相对 never buy 的差距是怎么随时间累积的。只描述看到的。
\item 如果想拆开表 \ref{tab:policies} 的差距(2{,}984 天每天差值的总和是 $-10{,}018$ bp,平均 $-3.36$ bp/天;用舍入后的 $-3.36$ 乘 2{,}984 得 $-10{,}026$,差的 8 bp 是舍入):仓位不同的 879 天占了 $-9{,}932$ bp。其中 never buy 做多、exact 空仓的有 333 天,而这 333 天里有 307 天 exact 在当天开始时就已经是空仓。这一条可以不写。
\end{pts}
\end{note}

\begin{model}
Under the fixed timer 18.2\% of cars were delayed (standard deviation 2.1 points across 30 runs); under the reactive signal 11.5\% were (1.8 points). In heavy rain the two controllers differed by less than the run-to-run spread, so we do not rank them there.
\end{model}

\subsection{Predicted and realized gain}

\begin{table}[t]
\caption{\capnote{一句说各列是什么(模型预测的好处、实际的好处、两者之差);一句说两边怎么保证可比(照精确解做、立刻停止,这两边每天都从 never buy 的昨日仓位出发);一句说最后一列是什么:对``实际减预测''这个差算了 100 个区间(5 种块长 $\times$ 20 个种子),这一列数的是其中不含 0 的个数。}}
\label{tab:gain}
\centering\small
\tab{tables/tab2_predicted_realized.tex}
\end{table}

\begin{note}[第 1 段,约 70 词:直接回答问题二]
\begin{pts}
\item $\lambda=1$:模型预测 $+0.84$ bp/天;实际 $-2.70$,区间 $[-4.6,\,-1.0]$;两者之差 $-3.54$,区间 $[-5.5,\,-1.8]$。
\item $\lambda=5$:预测 $+0.06$;实际 $-0.32$,区间 $[-1.5,\,+0.6]$;没有结论。
\item 稳健性,一句:100 个区间在 $\lambda=1$ 时全部不含 0,在 $\lambda=5$ 时全部包含 0。
\end{pts}
\end{note}

\begin{note}[第 2 段,约 90 词:亏在哪些日子($\lambda=1$)。只描述,不说原因]
\begin{pts}
\item 这一段拆的是表 \ref{tab:gain} 的 $-2.70$ bp/天:2{,}984 天合计 $-8{,}065$ bp(这是没有舍入的每日数值之和;$-2.70 \times 2{,}984$ 得 $-8{,}057$,差别是舍入)。不要和表 \ref{tab:policies} 的 $-10{,}018$ bp 混在一起。
\item 两边仓位不同的只有 374 天,这些天的仓位差别带来 $-7{,}770$ bp。
\item 买预测付的价格合计只有 295 bp。所以亏的主要不是买预测的钱,而是买了以后持有了不同的仓位。
\item 照精确解做的那一边,做多的日子略少:71.2\% 对 73.5\%。这是表 \ref{tab:gain} 的比较(每天从同一个仓位出发)。表 \ref{tab:policies} 的 long 列顺序相反(exact 77.88\%,never buy 73.53\%),因为那里每个策略沿着自己的路径走。两个都提的话,要说明这一点。
\item 图 \ref{fig:day}(a) 的那一天可以当例子,但要说明它是挑出来的:买了两个预测以后从做多转为空仓,第二天价格涨了 45.4 bp;never buy 那天的 net utility 是 $+45.23$ bp,exact 是 $-1.20$ bp。
\end{pts}
\end{note}

\begin{table}[t]
\caption{\capnote{一句说两列是哪两个时期;一句说 partial correlation 在这里是什么意思(在已有免费预测的情况下,付费预测还和 $r$ 有多少关系);一句说括号里是标准误、怎么算的(block bootstrap,每块 21 天)。}}
\label{tab:periods}
\centering\small
\tab{tables/tab3_fit_vs_test.tex}
\end{table}

\begin{note}[第 3 段,约 70 词:两个时期各是什么样,见表 3。只描述,解释留到第 5 节]
\begin{pts}
\item 每个数字都带标准误一起报:A 的 partial correlation 从 $+0.030$(0.025)到 $-0.043$(0.047);B 从 $+0.013$(0.032)到 $+0.104$(0.059);$r$ 的均值从 2.75(1.65)到 5.62(1.73)bp。
\item 这三项的变化都只有 1.2 到 1.4 个标准误。算法:变化除以 $\sqrt{\mathrm{se}_{\mathrm{fit}}^{2}+\mathrm{se}_{\mathrm{test}}^{2}}$,这样算是把两个时期当作相互独立。表里另外两行:$r$ 的标准差是 1.3 个标准误;$y_0$ 和 $r$ 的相关系数是 1.9 个标准误,不在这个范围里,要单独描述。按你们自己的读数规则,这是``分不开'',不是``变了''。不要写 A 的预测能力反转了。
\item 在拟合期,两个 partial correlation 的区间就已经包含 0。
\item 稳健性:A 在两个时期都是 100 个区间里 0 个不含 0。B 在测试期是 27 个,而且全部出现在最长的两种块长上,所以也不要写成``显著''。免费预测 $y_0$ 在测试期是 100 个里 100 个不含 0。
\end{pts}
\end{note}

% ---------------------------------------------------------------------
\section{Discussion and Limitations}
% ---------------------------------------------------------------------

\begin{note}[第 1 段,约 80 词:两个回答合起来说明什么。评分表:结论要从结果里推得出来]
\begin{pts}
\item 这一段的结论由你们判断。可以用的事实:
\item 问题一:仓位在所有价格下都相同。但只有两个预测时,两种做法只可能在第一个决策上不同,所以这个结果对``预测更多时贪心够不够好''说明不了多少。
\item 问题二:$\lambda=1$ 时模型预测的好处没有出现。$\lambda=5$ 没有结论,可以用阈值解释为什么:做多的阈值约 9.2 bp,never buy 有 98.4\% 的日子空仓,照精确解做的那一边和它仓位不同的只有 73 天(这是表 \ref{tab:gain} 的比较;按表 \ref{tab:policies} 各走各路的比较是 81 天)。
\item 参考行的含义。直接的比较是 exact 减 always long:$\lambda=1$ 时 $-3.34$ bp/天,区间 $[-5.8,\,-1.0]$。所以这里的损失也是相对``一直持有''的损失。(never buy 对 always long:区间包含 0,没有结论。不要从``分不开''去推。)
\item 一种可能的读法:这里出问题的不是``只看一步''。同意才用,并且到此为止,不再往前推。
\end{pts}
\end{note}

\begin{note}[第 2 段,约 80 词:表 3 和阈值放在一起能说什么。用 ``is consistent with'',不用 ``is caused by'']
\begin{pts}
\item agent 的模型依赖的关系,在拟合期就很弱,而且估计得不精确(表 \ref{tab:periods})。
\item $\lambda=1$ 时,拟合的均值 2.75 bp 刚好高于``入场''阈值 2.64 bp,比``继续做多''的阈值 0.64 bp 高约 2 bp。测试期里,买之前的 belief 均值有 38.6\% 的日子落在这两个阈值之间。
\item 所以一个小的向下更新就能让做多的 agent 离场,而重新入场需要更高的均值。写的时候说明这个说法和哪一种解释一致,不要写成已经证明的原因。
\item 注意表 \ref{tab:policies}:沿自己的路径走时,exact 做多的日子更多,不是更少(77.88\% 对 73.53\%)。所以不要写``买了以后做多更少''这种一般性的话。损失在 never buy 做多而 exact 没有做多的日子上(never buy 做多、exact 空仓的 333 天,合计 $-6{,}761$ bp)。
\item \citet{smith2006optimizer} 讲的``按估计值挑选会事后失望''和问题二直接相关,读了再决定引不引。
\end{pts}
\end{note}

\begin{note}[第 3 段,约 70 词:没有分开的几种可能。这一段决定论文是否诚实]
\begin{pts}
\item 规划深度:检验过了,在这里没有差别。
\item 另外三种没有分开:拟合模型的抽样误差;两个时期之间数据变了;模型族不对(高斯、线性)。
\item 对应课本的假设:抽样误差和模型族,对应``模型是给定的''(第 20 章开头,第 407 页:``given models of the transitions, observations, and rewards'';第 III 部分的引言说明这个假设在那里才放松)。数据随时间变,对应 7.1 节(第 133 页)的 stationary 假设。
\item status update 里写过原因是 non-stationary skill。这个说法没有被检验(见表 \ref{tab:periods} 的标准误)。用一句话如实更正。
\end{pts}
\end{note}

\begin{note}[第 4 段,约 70 词:局限和假设。评分表原文要求有]
\begin{pts}
\item 一个资产、一次时期划分。
\item 只有两个付费预测(所以问题一只检验了第一个决策);价格是人为设定的。
\item 规划器不跨天往前看。
\item belief 只拟合一次,而且是高斯的。
\item 研究问题是看到结果以后确定的。
\item 可选的一句,篇幅够并且重跑过才写:10 月的模拟器在 14{,}175 个设定里,贪心在 865 个(6.10\%)里比精确解差。它的参数和 SPY 实验不同(标准化单位下 $\lambda s = 0.05$、$\tau/s = 0.05$,即 $s=1\%$ 时 $\lambda=5$、$\tau=5$ bp;昨日空仓;belief 均值从交易边界到它上方一个标准差),必须带着参数一起写,不要和 SPY 的数字并排比较。这些数字的出处是 \texttt{analysis/a5\_output.txt}。
\end{pts}
\end{note}

\begin{note}[第 5 段,约 50 词:未来工作。评分表原文要求有]
\begin{pts}
\item 配对的对照实验:从拟合的模型里生成数据,再重新估计、重新回放。如果那里预测和实际一致,就能排除代码错误和一部分解释。
\item 把模型参数的不确定性也放进 belief。课本 16.4--16.5 节是这个想法,但那里处理的是离散的转移概率,不是高斯的均值和协方差,所以只当作想法来引。
\item 课本 14.3 节(第 288--289 页)讲策略对模型误差的稳健性,可以引一句,不要展开:这一章 12 月 3 日才讲。
\end{pts}
\end{note}

\begin{model}
We tested one intersection with simulated arrivals. We did not test whether the reactive signal still helps when the arrival rate changes, so our result should not be read as a claim about real traffic.
\end{model}

% ---------------------------------------------------------------------
\section{Conclusion}
% ---------------------------------------------------------------------

\begin{note}[约 70 词,五句左右]
\begin{pts}
\item 一句重复问题,两句给出两个回答(各带一个数字),一句说这对读课本 6.6 节的人意味着什么。
\item 最后一句点出最主要的局限和下一步实验。评分表写的是 ``The conclusion also includes discussion of limitations, assumptions, and future work'',所以结论里要有这一句。
\item 不出现前面没有的新内容,不出现新数字。
\end{pts}
\end{note}

% ---------------------------------------------------------------------
\section*{Contributions}
% ---------------------------------------------------------------------

\begin{note}[课程要求:结论之后写每个人做了什么]
两三句,写事实:每个人实际做了什么(想法、推导、审查和测试代码、跑实验、写哪几节)。要和下面的 AI 声明一致:由 AI 写的代码,这里就写谁审查、谁测试、谁做的决定,不要写成某人写的。
\end{note}

% ---------------------------------------------------------------------
\section*{Use of AI Tools}
% ---------------------------------------------------------------------

\begin{note}[课程要求:写明用了哪些 AI 工具、各自用来做什么。必须如实、写全]
下面是事实清单,句子你们自己写:
\begin{pts}
\item 工具:Claude Code(Anthropic 的 Claude 模型)。如果还用过别的(翻译软件、语法检查、其他聊天机器人),一并列出。
\item 代码和实验:写和调试代码与测试;在 GitHub Actions 上跑实验。
\item 分析:设计并运行了追加的分析(always long 和 always flat 两个参考行;块长和种子的稳健性检查;按仓位拆分损失);核对数字。
\item 图表和版式:设计并生成图和表;挑选了图 \ref{fig:day} 的示例日;这份带提示的提纲(中文版和英文版)、LaTeX 版式和排好的公式。
\item 研究问题和解读:AI 提出并比较了几种研究问题的提法,推荐了现在这两个问题;还提出了对结果的读法和``不要声称什么''的清单。采用了哪些、怎么改的,是你们的判断,如实写。
\item 文献:查找文献并核实条目。还有:解释概念、制定计划。
\item 正文由你们自己写。写完以后让 AI 核对了数字、逻辑和引用(只指出问题,不改写),这一条也写上。
\item status update 里的声明是 ``we used Claude Code to help write code; the experiment design, analysis, and all else are our own work''。实际用得比这句话多,这一次要写全。
\end{pts}
\end{note}

% ---------------------------------------------------------------------
%  References: every cited entry in refs.bib was checked on 2026-11-22
%  (Crossref, jmlr.org or the local PDF). One open item: the page range
%  of Politis and Romano (1992) comes from a secondary source.
% ---------------------------------------------------------------------
\bibliographystyle{unsrtnat}
\bibliography{refs}

\end{document}
